\subsection{Spectrum of an element in a normed algebra}

In this section, we assume that K = C.

THEOREM 1. --- Let A be a unital Banach algebra and let \(x \in \mathrm{A}\) .

\begin{enumerate}
\def\labelenumi{\alph{enumi})}
\item
  The set \(\mathrm{Sp}_{\mathrm{A}}(x)\) is a compact subset of C;
\item
  The spectral radius \(\varrho(x)\) is the radius of the smallest closed disk with center 0 in \(\mathbf{C}\) which contains \(\mathrm{Sp}_{\mathrm{A}}(x)\) ;
\item
  The resolvent \(\lambda \mapsto \mathrm{R}(x,\lambda) = (\lambda -x)^{-1}\) of \(x\) is holomorphic in \(\mathbf{C} - \mathrm{Sp}_{\mathrm{A}}(x)\) and zero at infinity. Moreover, for every integer \(k\geqslant 0\) , we have the
following formula:

\[
\frac {\partial^ {k}}{\partial \lambda^ {k}} \mathrm{R} (x, \lambda) = (- 1) ^ {k} k! \mathrm{R} (x, \lambda) ^ {k + 1};
\]

\item
  For every complex number \(\lambda\) such that \(|\lambda| > 1 / \varrho(x)\) , one has

\[
\mathrm{R} (x, \lambda) = \sum_ {n = 0} ^ {+ \infty} \lambda^ {- n - 1} x ^ {n}.
\]

\end{enumerate}

The complement of the spectrum \(\mathrm{Sp}_{\mathrm{A}}(x)\) is the inverse image of the group G of invertible elements of A under the continuous map \(\lambda \mapsto x - \lambda\) from C to A; by Proposition 3, b) of I, p.~23, the spectrum \(\mathrm{Sp}_{\mathrm{A}}(x)\) is a closed subset of C. Furthermore, let \(\lambda \in C\) such that \(|\lambda| > \varrho(x)\) . We have \(\lambda - x = \lambda(1 - \lambda^{-1}x)\) . Since \(\varrho(\lambda^{-1}x) = |\lambda|^{-1}\varrho(x) < 1\) , the element \(1 - \lambda^{-1}x\) , hence also the element \(\lambda - x\) , is invertible and

\[
\mathrm{R} (x, \lambda) = (\lambda - x) ^ {- 1} = \sum_ {n = 0} ^ {\infty} \lambda^ {- n - 1} x ^ {n}\tag{5}
\]

(I, p.~22, Prop. 2). In particular, \(\lambda \notin \mathrm{Sp}_{\mathrm{A}}(x)\) . This proves that \(\mathrm{Sp}_{\mathrm{A}}(x)\) is contained in the disk with center 0 and radius \(\varrho(x)\) . Consequently, \(\mathrm{Sp}_{\mathrm{A}}(x)\) is compact. This formula (5) also shows that the resolvent of \(x\) is defined and holomorphic in the complement of the closed disk \(\Delta_{\varrho(x)}\) of center 0 and radius \(\varrho(x)\) and tends to 0 at infinity.

Let \(\lambda_0\in \mathbf{C} - \mathrm{Sp}_{\mathrm{A}}(x)\) . Set \(y = \lambda_0 - x\) . Let \(\mu \in \mathbf{C}\) such that \(|\lambda_0 - \mu | < \| y^{-1}\|^{-1}\) . We have

\[
\mu - x = y - (\lambda_ {0} - \mu) = y (1 - (\lambda_ {0} - \mu) y ^ {- 1}),
\]

hence \(\mu - x\) is invertible and has as its inverse

\[
(\mu - x) ^ {- 1} = y ^ {- 1} \sum_ {n = 0} ^ {\infty} (\lambda_ {0} - \mu) ^ {n} y ^ {- n}\tag{6}
\]

By prop. 2 of I, p.~22. Hence the resolvent of \(x\) is defined and holomorphic in the open disk centered at \(\lambda_0\) and of radius \(\| y^{-1} \|^{-1}\) . Consequently, the resolvent of \(x\) is a holomorphic map from \(\mathbf{C} - \mathrm{Sp}_{\mathrm{A}}(x)\) into A.

Formula (1) of I, p.~4 implies \(\frac{\partial}{\partial\lambda}\mathrm{R}(x,\lambda) = -\mathrm{R}(x,\lambda)^2\) , whence, by induction on \(k\) ,

\[
\frac {\partial^ {k}}{\partial \lambda^ {k}} \mathrm{R} (x, \lambda) = (- 1) ^ {k} k! \mathrm{R} (x, \lambda) ^ {k + 1}.
\]

Let \(a > 0\) be a real number such that \(\mathrm{Sp}_{\mathrm{A}}(x)\) is contained in the closed disk \(\Delta_{a}\) of center 0 and radius \(a\) . The function \(\lambda \mapsto (\lambda^{-1} - x)^{-1}\) is then defined and holomorphic for \(0 < |\lambda| < a^{-1}\) and tends to 0 when \(\lambda\) tends to 0. The unique continuous function on the open disk centered at 0 with radius \(a^{-1}\) which extends this holomorphic function is then holomorphic (VAR, R1, 3.3.9), hence the radius of convergence of the series (5) defining it is \(\geqslant a^{-1}\) (VAR, R1, 3.2.9). By prop. 2 of I, p.~22, we therefore have \(a \geqslant \varrho(x)\) .

Remarks. --- 1) The spectrum of an element in a unital Banach algebra can be any arbitrary nonempty compact subset F of C (cf. Example 3 of I, p. 17; for A = C(F; C) and f ∈ A the canonical inclusion of F into C, one has \(\mathrm{Sp}_{\mathrm{A}}(f) = F\)).

\begin{enumerate}
\def\labelenumi{\arabic{enumi})}
\setcounter{enumi}{1}
\tightlist
\item
  Let A be a unital Banach algebra, and let \(x \in A\) . By Theorem 1 of I, p.~24, \(\mathbf{C} - \mathrm{Sp}_{\mathrm{A}}(x)\) is an open subset of C, hence is locally connected. Therefore the connected components of \(\mathbf{C} - \mathrm{Sp}_{\mathrm{A}}(x)\) are open. By Thm. 1, one of these connected components contains the set of \(\lambda \in C\) such that \(|\lambda| > \varrho(x)\) ; all the other connected components are therefore bounded.
\end{enumerate}

COROLLARY 1. --- Let A be a nonzero unital normed algebra. For every \(x \in \mathrm{A}\) the spectrum \(\mathrm{Sp}_{\mathrm{A}}(x)\) is nonempty.

Suppose first that A is complete. If one had \(\mathrm{Sp}(x)=\varnothing\) , the resolvent of x would be holomorphic in C and vanish at infinity, hence identically zero (VAR, R., 3.3.6, p.~29). Since \(\mathrm{R}(x,\lambda)=(\lambda-x)^{-1}\) is invertible, it would follow that 1=0 and hence that \(A=\{0\}\) .

In the general case, let \(\widehat{\mathrm{A}}\) be the completed algebra of A; the relation \(\mathrm{Sp}_{\mathrm{A}}(x) = \varnothing\) would entail \(\mathrm{Sp}_{\widehat{\mathrm{A}}} (x) = \varnothing\) , whence \(\widehat{\mathrm{A}} = \{0\}\) and \(\mathrm{A} = \{0\}\) .

COROLLARY 2 (Gelfand–Mazur theorem). --- Let A be a normed algebra over C. If A is a field, then A = C · 1.

If \(x \in A\) , there exists \(\lambda \in \mathbf{C}\) such that \(x - \lambda\) is non-invertible (cor. 1), hence \(x - \lambda = 0\) and \(x \in \mathbf{C} \cdot 1\) .

COROLLARY 3. --- Let A be a unital Banach algebra and let x be an invertible element of A such that \(\|x\|=\|x^{-1}\|=1\) . Then \(\mathrm{Sp}(x)\subset\mathbf{U}\) .

Let \(\Delta\) be the disk with center 0 and radius 1 in \(\mathbf{C}\) . By th. 1 b) and the fact that \(\varrho(x) \leqslant \|x\|\) , one has \(\operatorname{Sp}(x) \subset \Delta\) . Likewise, \(\operatorname{Sp}(x)^{-1} = \operatorname{Sp}(x^{-1}) \subset \Delta\) , whence the corollary (cf.~I, p.~2, remark 4).

COROLLARY 4. --- Let E be a complex Banach space, \(\mathcal{L}(\mathrm{E})\) the Banach algebra of continuous endomorphisms of E and A a nonzero subalgebra of \(\mathcal{L}(\mathrm{E})\) such that E is a simple A-pseudomodule (A, II, p.~176, Appendix).

\begin{enumerate}
\def\labelenumi{\alph{enumi})}
\item
  Let u be an endomorphism of E, not necessarily continuous, that commutes with A. Then u is a homothety.
\item
  Let \(u\) be an endomorphism of \(\mathbf{E}\) , not necessarily continuous. For every integer \(n \geqslant 1\) and for every \((\xi_1, \ldots, \xi_n) \in \mathrm{E}^n\) , there exists \(v \in A\) such that

\[
(v (\xi_ {1}), \dots , v (\xi_ {n})) = (u (\xi_ {1}), \dots , u (\xi_ {n})).
\]

\end{enumerate}

Let us show \(a\) ). Let \(\widetilde{\mathrm{A}}\) be the algebra obtained from A by adjoining an identity element. Since E is a simple A-pseudomodule, it is a simple \(\widetilde{\mathrm{A}}\)-module.

Let B be the commutant of \(\tilde{\mathsf{A}}\) in the ring of endomorphisms of the C-vector space E. The algebra B contains 1 and is the algebra of endomorphisms of the \(\tilde{\mathsf{A}}\)-module E. Since E is a simple \(\tilde{\mathsf{A}}\)-module, Schur's lemma (A, VIII, p.~43, corollary) shows that B is a field.

Let \(\xi_0\in \mathrm{E}\) be such that \(\mathrm{A}\xi_0\neq\{0\}\) . We therefore have \(\mathrm{A}\xi_0 = E\) . For every \(u\in \mathrm{B}\) , let \(\mathrm{A}_u\) be the set of \(v\in \mathrm{A}\) such that \(v(\xi_0) = u(\xi_0)\) . This set is nonempty, since \(\mathrm{A}\xi_0 = E\) . We then set

\[
\| u \| _ {\mathrm{B}} = \inf _ {v \in \mathrm{A} _ {u}} \| v \|.
\]

The map \(u \mapsto \|u\|_{\mathrm{B}}\) is a seminorm on B.

Let us show that this map is a norm. Let u be a nonzero element of B. For every \(v \in A_{u}\) , one has \(\|v\| \geqslant \|v(\xi_{0})\|/\|\xi_{0}\| = \|u(\xi_{0})\|/\|\xi_{0}\|\) , so that \(\|u\|_{B} \geqslant \|u(\xi_{0})\|/\|\xi_{0}\|\) . It therefore suffices to show that \(u(\xi_{0}) \neq 0\) . Let \(\xi_{1} \in E\) be such that \(u(\xi_{1}) \neq 0\) . Since \(A\xi_{0} = E\) , there exists \(w \in A\) such that \(\xi_{1} = w(\xi_{0})\) . Then, \(wu(\xi_{0}) = uw(\xi_{0}) = u(\xi_{1}) \neq 0\) , hence \(u(\xi_{0}) \neq 0\) .

On the other hand, let \(u\) and \(u'\) be elements of B. For every \(\varepsilon > 0\) , there exists \(v, v' \in A\) such that \(v(\xi_0) = u(\xi_0), v'(\xi_0) = u'(\xi_0)\) and \(\|v\| \leqslant \|u\|_{\mathrm{B}} + \varepsilon\) , \(\|v'\| \leqslant \|u'\|_{\mathrm{B}} + \varepsilon\) . Then we have \(vv'(\xi_0) = vu'(\xi_0) = u'v(\xi_0) = u'u(\xi_0)\) , whence

\[
\| u ^ {\prime} u \| _ {\mathrm{B}} \leqslant \| v ^ {\prime} v \| \leqslant \| v \| \| v ^ {\prime} \| \leqslant (\| u \| _ {\mathrm{B}} + \varepsilon) (\| u ^ {\prime} \| _ {\mathrm{B}} + \varepsilon),
\]

and finally \(\|u'u\|_{B} \leqslant \|u\|_{B} \|u'\|_{B}\) . This shows that B, equipped with the norm \(u \mapsto \|u\|_{B}\) , is a normed algebra. Since it is a field, Corollary 2 implies that \(B = C \cdot 1\) , which is the desired conclusion.

We prove \(b\) ). By \(a\) ), the commutant of A in \(\mathrm{End}_{\mathbf{C}}(\mathrm{E})\) is reduced to the homotheties of E. Its bicommutant is therefore \(\mathrm{End}_{\mathbf{C}}(\mathrm{E})\) . Assertion \(b\) ) therefore follows from the Jacobson density theorem (Theorem 1 of A, VIII, p.~434).

COROLLARY 5. --- Let A be a Banach algebra and \(x \in A\) .

\begin{enumerate}
\def\labelenumi{\alph{enumi})}
\item
  \(\mathrm{Sp}^{\prime}(x)\) is a compact subset of C;
\item
  The spectral radius \(\varrho(x)\) is the radius of the smallest closed disk centered at 0 in C that contains \(\mathrm{Sp}'(x)\) ;
\item
  For \(x\) to be quasi-nilpotent, it is necessary and sufficient that one has \(\mathrm{Sp}'(x) = \{0\}\) .
\end{enumerate}

The assertions \(a)\) and \(b)\) follow from th. 1 by considering the Banach algebra deduced from A by adjoining an identity element. The assertion \(c)\) follows from \(b)\) .

